\centerline{\bf \Huge Региональная тренировка I}

\begin{flushright}
        \scriptsize{Каждый ноль за задачу добавляет тыщу задач в домашку}
\end{flushright}

\problem{1}
Попарно различные натуральные числа $a, b$ и $c$ таковы, что числа $\frac{1}{a}, \frac{1}{b}$ и $\frac{1}{c}$ образуют возрастающую арифметическую прогрессию. Докажите, что наибольший общий делитель чисел $a$ и $b$ больше 1.

\problem{2}
Хорды $A C$ и $B D$ окружности пересекаются в точке $P$. Перпендикуляры к $A C$ и $B D$, проходящие через точки $C$ и $D$ соответственно, пересекаются в точке $Q$. Докажите, что прямые $A B$ и $P Q$ перпендикулярны.

\problem{3}
Дано вещественное число $x>1$. Известно, что $\left[x^2\right],\left[x^3\right],\left[x^4\right]-$ полные квадраты. Докажите, что $[x]$ так же является полным квадратом.

\problem{4}
На плоскости расположены $n$ квадратов $2 \times 2$ со сторонами, параллельными координатным осям. Ни один из центров этих квадратов не содержится ни в каком другом квадрате (в том числе на границе). Прямоугольник П со сторонами, параллельными координатным осям, содержит все эти квадраты. Докажите, что периметр П не меньше $4 \sqrt{n}$.

\problem{5}
Пусть $x, y$ - целые числа. Последовательность $\left\{a_n\right\}$ определяется так:

$$
a_0=a_1=0, a_{n+2}=x a_{n+1}+y a_n+1
$$


для целых неотрицательных $n$. Докажите, что для простых $p$ верно, что НОД$\left(a_p, a_{p+1})\right.$ либо равен 1 , либо больше $\sqrt{p}$.